\section{Conclusion}
We have shown that metric conditioning and spectral-inversion stability can separate
qualitatively near a pseudo-Hermitian exceptional point.  In a two-level family the optimized
metric certificate is sharp, whereas a coupled three-level family approaches a defective
nonzero eigenvalue while remaining uniformly invertible in Euclidean norm.  For that family,
the minimum condition number over all positive compatible metrics obeys the sharp law
\[
 \kappa_*(A_\varepsilon)=\frac{8}{\varepsilon^2}\,[1+o(1)],
 \qquad \varepsilon\downarrow0,
\]
and an explicit metric attains the leading constant.  Consequently the best metric-induced
Euclidean invertibility certificate vanishes as $\varepsilon/(2\sqrt2)$ even though the true
distance to singularity stays positive.  The same separation persists in a genuinely coupled
growing-dimensional chain, while a standard Hermitian dilation retains uniformly bounded
singular-value conditioning.

The resulting message is precise: singularity of the physical metric records the collapse of
eigenvector geometry required to maintain quasi-Hermiticity, but it need not signal instability
of the inverse problem in the reference norm.  Distinguishing these two effects is essential
when metric conditioning is used to interpret spectral stability or computational cost near
exceptional points.

\section*{Data availability}
All numerical data supporting the figures are generated by the reproducibility scripts available
in the public GitHub repository
\url{https://github.com/divinelybless/Metric-Conditioning-and-Stable-Spectral-Inversion-near-Exceptional-Points-in-Pseudo-Hermitian-System}.
No external research dataset is required to verify the analytical results.

\section*{Funding}
This research received no external funding.

\section*{Conflict of interest}
The authors declare no competing interests.

\begin{thebibliography}{99}
\bibitem{Mostafazadeh2002I}
A. Mostafazadeh, ``Pseudo-Hermiticity versus PT symmetry: The necessary condition for the reality of the spectrum of a non-Hermitian Hamiltonian,'' \emph{J. Math. Phys.} \textbf{43}, 205--214 (2002). \href{https://doi.org/10.1063/1.1418246}{doi:10.1063/1.1418246}.
\bibitem{Mostafazadeh2002II}
A. Mostafazadeh, ``Pseudo-Hermiticity versus PT symmetry. II. A complete characterization of non-Hermitian Hamiltonians with a real spectrum,'' \emph{J. Math. Phys.} \textbf{43}, 2814--2816 (2002). \href{https://doi.org/10.1063/1.1461427}{doi:10.1063/1.1461427}.
\bibitem{Mostafazadeh2002III}
A. Mostafazadeh, ``Pseudo-Hermiticity versus PT symmetry. III. Equivalence of pseudo-Hermiticity and the presence of antilinear symmetries,'' \emph{J. Math. Phys.} \textbf{43}, 3944--3951 (2002). \href{https://doi.org/10.1063/1.1489072}{doi:10.1063/1.1489072}.
\bibitem{Krejcirik2018}
D. Krej\v{c}i\v{r}\'{i}k, V. Lotoreichik, and M. Znojil, ``The minimally anisotropic metric operator in quasi-Hermitian quantum mechanics,'' \emph{Proc. R. Soc. A} \textbf{474}, 20180264 (2018). \href{https://doi.org/10.1098/rspa.2018.0264}{doi:10.1098/rspa.2018.0264}.
\bibitem{Trefethen2005}
L. N. Trefethen and M. Embree, \emph{Spectra and Pseudospectra: The Behavior of Nonnormal Matrices and Operators} (Princeton University Press, 2005).
\bibitem{Gilyen2019}
A. Gily\'{e}n, Y. Su, G. H. Low, and N. Wiebe, ``Quantum singular value transformation and beyond: Exponential improvements for quantum matrix arithmetics,'' in \emph{Proceedings of STOC 2019}, pp. 193--204 (2019). \href{https://doi.org/10.1145/3313276.3316366}{doi:10.1145/3313276.3316366}.
\bibitem{Berry2026}
M. E. S. Morales, L. Pira, P. Schleich, K. Koor, P. C. S. Costa, D. An, A. Aspuru-Guzik, L. Lin, P. Rebentrost, and D. W. Berry, ``Quantum linear system solvers: A survey of algorithms and applications,'' \emph{Rev. Mod. Phys.} \textbf{98}, 025005 (2026). \href{https://doi.org/10.1103/x6gh-d8gh}{doi:10.1103/x6gh-d8gh}.
\bibitem{Schomerus2024}
H. Schomerus, ``Eigenvalue sensitivity from eigenstate geometry near and beyond arbitrary-order exceptional points,'' \emph{Phys. Rev. Research} \textbf{6}, 013044 (2024). \href{https://doi.org/10.1103/PhysRevResearch.6.013044}{doi:10.1103/PhysRevResearch.6.013044}.
\end{thebibliography}

\end{document}
